% Generated draft: requires mathematical and source review.
\documentclass[11pt]{article}
\usepackage[margin=25mm]{geometry}
\usepackage{fontspec}
\setmainfont{Latin Modern Roman}
\usepackage{amsmath,amssymb,mathtools,unicode-math}
\setmathfont{Latin Modern Math}
\usepackage{enumitem,needspace,xurl,hyperref}
\setlength{\parindent}{0pt}
\setlength{\parskip}{5pt}
\setlength{\emergencystretch}{3em}
\newcommand{\N}{\mathbb{N}}
\newcommand{\Z}{\mathbb{Z}}
\newcommand{\Q}{\mathbb{Q}}
\newcommand{\R}{\mathbb{R}}
\newcommand{\C}{\mathbb{C}}
\newcommand{\F}{\mathbb{F}}
\newcommand{\A}{\mathbb{A}}
\newcommand{\PP}{\mathbb{P}}
\newcommand{\E}{\mathbb{E}}
\newcommand{\eps}{\varepsilon}
\newcommand{\dd}{\,\mathrm{d}}
\begin{document}
\section*{UnsolvedMath ID: 30005495}
\textbf{Source problem code:} catalogue entry
\subsection{Definitions and mathematical statement}
A group \(G\) is finitely presented if there exist finite sets of generators and relations such that \(G\cong\langle g_1,\ldots,g_r\mid R_1,\ldots,R_s\rangle\).

For a conventional complex, boundary-free interpretation, a log-canonical singularity is a normal complex germ \((X,x)\) whose canonical divisor \(K_X\) is \(\Q\)-Cartier and whose discrepancies are at least \(-1\). More explicitly, for every log resolution \(f:Y\to X\), writing
\[
K_Y=f^*K_X+\sum_{E\text{ exceptional}}a(E,X)E,
\]
one requires \(a(E,X)\geq -1\) for every exceptional prime divisor \(E\). This is the log-canonical condition for the pair \((X,0)\).

The supplied question asks whether
\[
\forall G\text{ finitely presented},\quad
\exists (X,x)\text{ log-canonical with }\dim X=5
\quad\text{such that}\quad
\pi_1^{\mathrm{loc}}(X,x)\cong G.
\]
Here dimension is complex dimension under the conventional complex interpretation above. The notation \(\pi_1^{\mathrm{loc}}\) requires a source convention that is not specified in the supplied extract. In particular, two possible topological meanings are
\[
\pi_1(U\setminus\{x\},b)
\qquad\text{and}\qquad
\pi_1(U\cap X_{\mathrm{reg}},b),
\]
where \(U\) is a sufficiently small representative neighborhood of \(x\), \(X_{\mathrm{reg}}\) is the smooth locus, and \(b\in U\cap X_{\mathrm{reg}}\). The first is the punctured-neighborhood interpretation, equivalently the fundamental group of the link in the usual complex analytic setting. The second removes the entire singular locus. They coincide for an isolated singularity but need not coincide in general.

\emph{Unresolved qualifications.} The record imposes no isolated-singularity hypothesis and does not specify which local fundamental group is intended, whether the category is algebraic or analytic, or the ground field explicitly. The complex, boundary-free definition above is therefore a conditional interpretation, not a verified expansion of the source. The realization question must retain these unresolved conventions pending source review; no isolatedness assumption or boundary divisor is added.
\subsection{Short English statement}
The question asks whether every finitely presented group can be realized as the local fundamental group of a log-canonical singularity of dimension five. The meaning of local fundamental group, the geometric category, and the precise singularity conventions require confirmation. The supplied record labels the problem open; this draft does not independently verify that status.
\subsection{Sources}
\begin{itemize}
\item Supplied local UnsolvedMath registry record, permanent ID \(30005495\), in the Oberwolfach Reports Open Problems collection; exact source problem code \texttt{catalogue entry}. The supplied original wording is: ``Does any finitely presented group appear as the fundamental group of a 5-dimensional lc singularity?'' The cleaned wording specifies ``local fundamental group.''
\item \emph{Komplexe Analysis -- Differential and Algebraic methods in Kähler spaces} (2023), \emph{Oberwolfach Reports}, pp. 965--1030, DOI \texttt{10.4171/owr/2023/18}. The original record points to p. 1009. Supplied source link: \url{https://doi.org/10.4171/owr/2023/18}. No individual proposer is identified in the supplied record.
\item Supplied literature source labeled ``Low-dimensional lc singularities'': \url{https://arxiv.org/abs/2302.11790}.
\item Supplied literature source labeled ``OWR 2023 report'': \url{https://ems.press/content/serial-article-files/47013}.
\item Literature assessment dated exactly \texttt{2026-08-22}: the supplied curation labels the problem open and states that the five-dimensional log-canonical realizability question for arbitrary finitely presented groups remains unresolved, with current low-dimensional results not reaching this claim. This is supplied dated curation, not an independently verified literature assessment; the record itself requests independent verification of time-sensitive status and cited claims.
\item The supplied statement status is \texttt{corrected\_verified}, accompanied by a claim that the cleaned statement was checked against the official report and preserves its hypotheses, notation, and target. That verification claim is source metadata and has not been independently checked for this draft.
\end{itemize}
\end{document}
